\begin{proof}
\textbf{Theorem (Bernoulli’s inequality).}  
Let $x\in\mathbb R$ with $x\ge-1$ and let $n\in\mathbb N$ (i.e. $n\ge0$). Then
\[
(1+x)^{\,n}\;\ge\;1+n\,x .
\]

\medskip
\textbf{Base case $n=0$.}
\[
(1+x)^{0}=1=1+0\cdot x .
\]
Hence the inequality holds with equality for every admissible $x$.

\medskip
\textbf{Inductive step.}  
Assume that for some $n\ge0$ and for every $x\ge-1$
\begin{equation*}
(1+x)^{\,n}\;\ge\;1+n x .\tag{IH}
\end{equation*}
Because $x\ge-1$ we have $1+x\ge0$, so multiplication by the non‑negative factor $1+x$ preserves the inequality direction:
\begin{align*}
(1+x)^{\,n+1}
   &=(1+x)^{\,n}(1+x)\\
   &\ge (1+n x)(1+x). \tag{1}
\end{align*}
Expanding the right–hand side gives
\begin{align*}
(1+n x)(1+x)=1+(n+1)x+n x^{2}. \tag{2}
\end{align*}
Since $n\ge0$ and $x^{2}\ge0$, the term $n x^{2}$ is non‑negative. Dropping this non‑negative term yields the weaker (hence still valid) inequality
\begin{align*}
1+(n+1)x+n x^{2}\;\ge\;1+(n+1)x .\tag{3}
\end{align*}
Combining \eqref{1}–\eqref{3} we obtain
\[
(1+x)^{\,n+1}\;\ge\;1+(n+1)x ,
\]
which is precisely the desired statement for the exponent $n+1$.

\medskip
\textbf{Conclusion.}  
The base case $n=0$ is true, and the inductive step shows that truth for a given $n$ implies truth for $n+1$. By the principle of mathematical induction the inequality holds for every integer $n\ge0$ and every real $x\ge-1$.

\medskip
\textbf{Equality cases.}  
Equality can occur only in the steps where a non‑negative term was discarded.

\begin{itemize}
\item In the base case $n=0$ we have equality for all admissible $x$.
\item In the inductive step the term $n x^{2}$ was removed. Equality in that step requires $n x^{2}=0$, i.e.\ either $n=0$ (the base case) or $x=0$.
\end{itemize}
Consequently the original inequality becomes an equality precisely in the following situations
\[
\begin{cases}
n=0 &\text{(any }x\ge-1),\\[4pt]
n=1 &\text{(any }x\ge-1),\\[4pt]
x=0 &\text{(any }n\ge0).
\end{cases}
\]
(For $x=-1$ the inequality reads $0\ge 1-n$, and equality holds only when $n=1$, which is already included above.)  
Thus equality holds \emph{iff} $x=0$ or $n\le1$.\qed
\end{proof}